\newcommand{\DoNotLoadEpstopdf}{}
\pdfinfoomitdate=1
\pdftrailerid{}
\pdfsuppressptexinfo=-1
\documentclass[11pt,letterpaper]{article}
\usepackage[T1]{fontenc}
\usepackage{lmodern}
\usepackage{amsmath,amssymb,amsthm,mathtools}
\usepackage[margin=1in]{geometry}
\usepackage{booktabs,array,microtype}
\usepackage[hidelinks]{hyperref}
\hypersetup{pdftitle={Report 164 Exact generating function for pop stacked permutations},pdfauthor={},pdfsubject={Explicit EGF, poles, differential algebraicity and faster enumeration of A307030}}
\setlength{\parskip}{4pt}
\setlength{\parindent}{1.2em}
\setlength{\emergencystretch}{3em}
\setcounter{tocdepth}{1}
\numberwithin{equation}{section}
\newtheorem{theorem}{Theorem}[section]
\newtheorem{lemma}[theorem]{Lemma}
\newtheorem{proposition}[theorem]{Proposition}
\newtheorem{corollary}[theorem]{Corollary}
\theoremstyle{remark}\newtheorem{remark}[theorem]{Remark}
\newcommand{\dd}{\,\mathrm d}
\newcommand{\one}{\mathbf1}
\newcommand{\CC}{\mathbb C}
\newcommand{\lo}{\operatorname{lo}}
\newcommand{\hi}{\operatorname{hi}}
\newcommand{\spr}{\operatorname{spr}}
\newcommand{\Res}{\operatorname{Res}}
\newcommand{\ceil}[1]{\left\lceil#1\right\rceil}
\title{Exact generating function for pop stacked permutations\\\large Poles differential algebraicity and faster enumeration}
\author{Report 164}
\date{3 October 2026}
\begin{document}
\maketitle
\begin{abstract}
Let $p_n$ count permutations in the image of one pass through a pop stack,
with $p_0=1$. We derive an explicit exponential generating function
$P(z)=\sum_{n\geq0}p_nz^n/n!$ for OEIS A307030. Its branch-free expression
is a quotient of entire functions. A cumulative integral boundary problem
and an explicit half-angle cancellation establish the formula from the
run-overlap characterization, rather than from fitted coefficients.

Every genuine complex pole is simple and has an explicit residue.
There are infinitely many positive real poles, proving that $P$ is not
D-finite and $(p_n)$ is not P-recursive. Nevertheless $P$ is differentially
algebraic, with an order-at-most-two polynomial annihilator given by a
small resultant. An exact Riccati equation computes $p_0,\ldots,p_N$ in
$O(N^2)$ arithmetic operations and $O(N)$ scalar storage, improving the
$O(N^4)$ arithmetic-operation total-count algorithms in the cited papers.

A self-contained positive-operator argument isolates the unique dominant
pole $\rho$. We obtain its exact amplitude $C$, rationally certify narrow
intervals for both, and prove finite-radius pole expansions. Fixed-order
Stirling and rounding-safe inverse consequences follow. The numerical
pole estimates and leading constants are prior art; the contribution here
is the explicit formula and proofs. No effective numerical complex gap,
infinite pole-sum convergence, bit-complexity, optimality, or unrestricted
priority claim is made.
\end{abstract}
\newpage
\tableofcontents
\medskip
\noindent\textbf{Scope and reproduction.}
This is a standalone derivation. The overlap characterization and earlier
counting algorithms are attributed in Section~\ref{sec:sources}. A bounded
standard-library companion independently expands the entire quotient,
checks the integer recurrence, and certifies the constants with rational
inequalities. Its checks do not replace the proofs. All computed public
commands write deterministic JSON to standard output only. The source
package includes the report, reproducible PDF builder, fixtures, tests,
and checked output certificates. The existing Report 161 is unchanged.
\newpage

\section{The explicit generating function}\label{sec:main}
A pop-stack pass reverses each maximal decreasing run and concatenates
the resulting increasing blocks. A permutation in its image is
\emph{pop-stacked}. The first values, including the empty permutation, are
\[
1,1,1,3,11,49,263,1653,11877,95991,862047,8516221.
\]
The established characterization~\cite{abbhl,abh} says that consecutive
maximal ascending runs, with endpoints $(a,b)$ and $(c,d)$, satisfy
\begin{equation}\label{eq:overlap}
 a<d,\qquad c<b.
\end{equation}
A singleton has equal endpoints. Thus a singleton may sit strictly inside
the range of an adjacent run, but two singleton runs cannot be adjacent.
The second inequality enforces the descent between maximal runs; the
first is the extra image condition.

Put
\[
 Q(z)=4e^z-1,\qquad E(z)=e^{z/2},\qquad B(z)=2E(z)+1,
\]
and define the entire functions
\begin{align}
 H(z)&=\sum_{k\geq0}\frac{(-1)^kz^{2k}Q(z)^k}{16^k(2k)!},\label{eq:H}\\
 J_0(z)&=\frac z4\sum_{k\geq0}\frac{(-1)^kz^{2k}Q(z)^k}{16^k(2k+1)!}.
 \label{eq:J}
\end{align}
Factorial denominators give locally uniform convergence. Set
\begin{equation}\label{eq:ND}
 N(z)=e^{-z}Q(z)\bigl(H(z)+B(z)J_0(z)\bigr),\qquad
 D(z)=B(z)H(z)-Q(z)J_0(z).
\end{equation}
\begin{theorem}\label{thm:main}
The exponential generating function of pop-stacked permutations is
\begin{equation}\label{eq:entire}
 \boxed{P(z)=\frac{N(z)}{D(z)}}.
\end{equation}
This denotes a meromorphic function on $\CC$, with removable common zeros
filled in. Near zero, write $r=\sqrt{4e^z-1}$ with $r(0)=\sqrt3$; then
\begin{equation}\label{eq:trig}
 P(z)=e^{-z}r\,
 \frac{(2E+1)\sin(zr/4)+r\cos(zr/4)}
      {(2E+1)\cos(zr/4)-r\sin(zr/4)}.
\end{equation}
The global expression \eqref{eq:entire} uses no square-root branch.
The genuine pole set is exactly
\begin{equation}\label{eq:poleset}
 \{\zeta\in\CC:D(\zeta)=0,\ Q(\zeta)\ne0\}.
\end{equation}
Every member is simple, and
\begin{equation}\label{eq:residue}
 \Res(P,\zeta)=-\frac{2e^{-2\zeta}(4e^\zeta-1)}{\zeta+2}.
\end{equation}
There are infinitely many positive poles. The EGF is not D-finite, and
the counts are not P-recursive, although the EGF is differentially
algebraic over $\mathbb Q(z)$.
\end{theorem}
The proof occupies Sections~\ref{sec:model}--\ref{sec:algebraic}.
The exact recurrence and arithmetic complexity appear in
Section~\ref{sec:recurrence}; the leading asymptotic and constants appear
in Sections~\ref{sec:dominance}--\ref{sec:certificate}.

\section{Mixed run states and exact normalization}\label{sec:model}
Let
\[
 I=(0,1),\qquad \Delta=\{(a,b):0<a<b<1\},\qquad
 X=I\sqcup\Delta,\qquad \mu=\dd x\sqcup\dd a\dd b.
\]
A state $x\in I$ is a singleton with $\lo(x)=\hi(x)=x$; a state
$(a,b)\in\Delta$ has $\lo=a,\hi=b$. Define
\begin{align}
 K(s,t)&=\one\{\lo(s)<\hi(t),\ \lo(t)<\hi(s)\},\label{eq:K}\\
 w_z(x)&=z,\qquad w_z(a,b)=z^2e^{z(b-a)},\label{eq:weight}\\
 (A_zf)(s)&=\int_X K(s,t)w_z(t)f(t)\dd\mu(t).
 \label{eq:operator}
\end{align}
Work, for example, on $L^2(X,\mu)$. The kernel is bounded, $\mu(X)=3/2$,
and the weights are uniformly bounded on compact $z$-sets. Thus $A_z$
is compact, entire in operator norm, and $A_z=O(|z|)$ near zero.

Here is the factorial normalization. Represent $n$ distinct labels by
independent coordinates in $(0,1)$. Each of the $n!$ rank orders has
volume $1/n!$. A run of length $\ell\geq2$ with extreme coordinates
$a<b$ has ordered interior-coordinate volume
$(b-a)^{\ell-2}/(\ell-2)!$. Its summed weight is
\[
 \sum_{\ell\geq2}z^\ell\frac{(b-a)^{\ell-2}}{(\ell-2)!}
 =z^2e^{z(b-a)}.
\]
A singleton contributes $z\dd x$. The product of kernels enforces exactly
the adjacent-run conditions \eqref{eq:overlap}; there is neither an
extra run-order factor nor an extra factorial. For instance,
\[
 \int_Xw_z\dd\mu=z+z^2\int_0^1(1-t)e^{zt}\dd t=e^z-1,
\]
the EGF of the unique increasing permutation of each positive size.

Summing over the number of runs, with absolute convergence for small
$|z|$, now proves
\begin{equation}\label{eq:resolvent}
 h=(I-A_z)^{-1}\one,\qquad
 P(z)=1+\int_Xw_z h\dd\mu.
\end{equation}
The pairing is bilinear; no complex conjugate is inserted.

\section{Cumulative boundary problem}\label{sec:bvp}
For small $z$, define the endpoint cumulatives
\begin{align*}
 U(x)&=\int_{\hi(s)\leq x}w_z(s)h(s)\dd\mu(s),&
 V(x)&=\int_{\lo(s)<x}w_z(s)h(s)\dd\mu(s),\\
 T&=\int_Xw_z h\dd\mu=P-1.
\end{align*}
Endpoint equality sets have measure zero. Hence
\begin{equation}\label{eq:ends}
 U(0)=V(0)=0,\qquad U(1)=V(1)=T.
\end{equation}
Complementing the two nonoverlap regions in \eqref{eq:K} gives
\begin{equation}\label{eq:hUV}
 h(x)=1+V(x)-U(x),\qquad h(a,b)=1+V(b)-U(a).
\end{equation}
These identities first yield continuous representatives and then the
smoothness used below. Introduce
\[
 I_0(x)=\int_0^x e^{z(x-a)}U(a)\dd a,\qquad
 J_1(x)=\int_x^1 e^{z(b-x)}V(b)\dd b.
\]
Direct differentiation of the cumulatives gives
\begin{align}
 U'&=ze^{zx}(1+V)-zU-z^2I_0, & I_0'&=U+zI_0,\label{eq:firstU}\\
 V'&=zV+ze^{z(1-x)}(1-U)+z^2J_1, & J_1'&=-V-zJ_1.
 \label{eq:firstV}
\end{align}
For example, differentiating the upper endpoint of the interval part of
$U$ gives $z^2\int_0^xe^{z(x-a)}(1+V(x)-U(a))\dd a$; combining this with
the singleton term $z(1+V(x)-U(x))$ gives \eqref{eq:firstU}.
Differentiating again and substituting eliminates the integral terms:
\begin{equation}\label{eq:secondUV}
 U''=ze^{zx}V',\qquad V''=-ze^{z(1-x)}U'.
\end{equation}
Writing $u=U'$ and $v=V'$ yields
\begin{equation}\label{eq:uv}
 u'=ze^{zx}v,\qquad v'=-ze^{z(1-x)}u,
 \qquad u(0)=z,\quad v(1)=z.
\end{equation}
The final boundary condition uses $U(1)=V(1)$ in \eqref{eq:firstV}.
For $z\ne0$, elimination gives the constant-coefficient problem
\begin{equation}\label{eq:ode}
 u''-zu'+z^2e^zu=0,\qquad u(0)=z,\quad u'(1)=z^2e^z.
\end{equation}
At zero the homogeneous scalar problem is $u''=0,u(0)=u'(1)=0$;
its only solution is zero. Analytic fundamental matrices therefore give
local uniqueness for \eqref{eq:ode}. The first-order problem
\eqref{eq:uv} is likewise locally unique, including $z=0$.

\begin{lemma}[Sufficiency of the boundary problem]
The locally unique solution of \eqref{eq:uv}, integrated from zero,
reconstructs precisely the resolvent in \eqref{eq:resolvent}.
\end{lemma}
\begin{proof}
The reflected pair $(v(1-x),u(1-x))$ solves the same equations and boundary
data. Uniqueness gives $u(x)=v(1-x)$. Therefore the reconstructed
$U(x)=\int_0^xu$ and $V(x)=\int_0^xv$ satisfy $U(1)=V(1)$.
Define $I_0,J_1$ from these functions and put
\begin{align*}
 F_1&=u-ze^{zx}(1+V)+zU+z^2I_0,\\
 F_2&=v-zV-ze^{z(1-x)}(1-U)-z^2J_1.
\end{align*}
Calculation using \eqref{eq:uv} gives $F_1'=zF_1$ and $F_2'=-zF_2$.
Their boundary values are $F_1(0)=0$ and
$F_2(1)=z(U(1)-V(1))=0$. Both vanish. Consequently the endpoint
cumulatives of the $h$ defined by \eqref{eq:hUV} have derivatives
$U',V'$ and initial values zero; they are exactly $U,V$.
Complementing nonoverlap regions gives $h=1+A_zh$. The Neumann inverse
is unique near zero, completing the sufficiency proof.
\end{proof}
This step retains the endpoint condition that could otherwise be lost
by differentiating. No global uniqueness of the unreduced boundary
problem is asserted at exceptional complex parameters.

\section{Explicit solution and the canceled factor}\label{sec:cancellation}
Choose $r=\sqrt{4e^z-1}$ locally at zero and put $w=zr/2$. Solving
\eqref{eq:ode} gives
\begin{align}
 u(x)&=e^{zx/2}\bigl(z\cos(wx)+A\sin(wx)\bigr),\label{eq:usolution}\\
 \frac Az&=\frac{2E-\cos w+r\sin w}{r\cos w+\sin w}.
 \label{eq:A}
\end{align}
Integrate the differential equation and use $P=1+\int_0^1u$ to obtain
\begin{equation}\label{eq:integrated}
 P=\frac{zu(1)-z^2/2+Aw}{z^2E^2}.
\end{equation}
Substitution, with $\sin^2w+\cos^2w=1$, simplifies this to
\begin{equation}\label{eq:fullangle}
 P=e^{-z}r\frac{2E-\cos w+r\sin w}{r\cos w+\sin w}.
\end{equation}
All identities with apparent division by $z$ are first obtained for
small nonzero $z$ and then continued to zero.

To display the cancellation, set $c=\cos(w/2)$, $s=\sin(w/2)$ and
\[
 D_+=(2E+1)c-rs,\quad D_-=(2E-1)c+rs,\quad L=rc+(2E+1)s.
\]
The identity $r^2=(2E-1)(2E+1)$ gives, by multiplication,
\begin{align}
 r(r\cos w+\sin w)&=D_+D_-,\label{eq:factor1}\\
 r(2E-\cos w+r\sin w)&=D_-L.\label{eq:factor2}
\end{align}
As $D_-(0)=1$, the common factor cancels near zero. Substituting into
\eqref{eq:fullangle} proves \eqref{eq:trig}. In particular, zeros of the
full boundary determinant coming solely from $D_-$ are not poles of the
scalar counting EGF.

Locally $H=\cos(zr/4)$ and $J_0=\sin(zr/4)/r$, so \eqref{eq:trig}
is exactly \eqref{eq:entire}. Both $N,D$ are entire and $N(0)=D(0)=3$.
This proves the branch-free meromorphic continuation and $P(0)=1$.
It does not require a square root or inverse trigonometric function to
be defined globally.

\section{Positive poles and non D finiteness}\label{sec:positive}
For real $z\geq0$ put $r=\sqrt{4e^z-1}>0$ and
\begin{equation}\label{eq:phase}
 \alpha(z)=\frac12\arccos\!\left(\frac1{2e^{z/2}}\right),\qquad
 \phi(z)=\frac{zr}{4}+\alpha(z).
\end{equation}
The real inverse cosine takes values in $[0,\pi]$.
Since $\tan\alpha=r/(2E+1)$, equation \eqref{eq:trig} becomes
\begin{equation}\label{eq:tan}
 P(z)=e^{-z}r\tan\phi(z).
\end{equation}
Differentiation simplifies to
\begin{equation}\label{eq:phiprime}
 \phi'(z)=\frac{e^z(z+2)}{2r}>0.
\end{equation}
Moreover $\phi(0)=\pi/6$ and $\phi(z)\to\infty$. Thus for each
integer $k\geq0$ there is a unique positive $\rho_k$ satisfying
\begin{equation}\label{eq:realpoles}
 \phi(\rho_k)=\frac\pi2+k\pi,
\end{equation}
or equivalently
\[
 \frac{\rho_k\sqrt{4e^{\rho_k}-1}}2+
 \arccos\!\left(\frac1{2e^{\rho_k/2}}\right)=(2k+1)\pi.
\]
The derivative and prefactor in \eqref{eq:tan} are nonzero there, so all
these points are genuine simple poles. This describes positive poles
only, without ordering them among complex poles by modulus.

An analytic D-finite germ satisfies a nonzero linear differential
equation with polynomial coefficients. At every point where its leading
coefficient is nonzero, the standard local linear-ODE continuation
theorem excludes singularities. Its finite-plane singularities can thus
occur at only finitely many points. The infinitely many distinct poles
above prove that $P$ is not D-finite. If $(p_n)$ were P-recursive, then
$(p_n/n!)$ would be P-recursive and its ordinary generating function $P$
would be D-finite, a contradiction. The ordinary series
$\sum p_nz^n$, understood formally, is not D-finite either.

\section{Riccati equation and every complex pole}\label{sec:riccati}
On a local branch in \eqref{eq:tan}, let $a=e^{-z}r$. Then
\[
 \frac{a'}a=\frac{1-2e^z}{Q},\qquad
 a\phi'=\frac{z+2}{2},\qquad
 \frac{\phi'}a=\frac{e^{2z}(z+2)}{2Q}.
\]
Differentiating $P=a\tan\phi$ yields the exact identity
\begin{equation}\label{eq:riccati}
 2Q P'=2(1-2e^z)P+(z+2)Q+(z+2)e^{2z}P^2,
 \qquad P(0)=1.
\end{equation}
It holds globally by meromorphic continuation. Since $2Q(0)=6$, it is
also a locally unique analytic initial-value equation at zero.

\subsection{Removal of the degenerate parameters}
All zeros of $Q$ are $z_k=-\log4+2\pi ik$, and they are simple. At such
a point $E=(-1)^k/2$. If $k$ is even, $B=2$, $D=2$ and $P(z_k)=0$.
If $k$ is odd, locally
\[
 B=1-\sqrt{1+Q}=-Q/2+O(Q^2),\quad H=1+O(Q),\quad
 J_0=z/4+O(Q).
\]
Consequently
\[
 D=-Q(z+2)/4+O(Q^2),\qquad N=4Q+O(Q^2).
\]
As $z_k\ne-2$, the singularity is removable and
$P(z_k)=-16/(z_k+2)$. Thus no zero of $Q$ is a pole.

The other potentially degenerate coefficient for a Riccati pole balance
is $z=-2$. Put $s=\sqrt{1-4/e^2}$, so $0<s<1$. Direct evaluation gives
\[
 D(-2)=(1+2/e)\cosh(s/2)-s\sinh(s/2)>0,
\]
because $1+2/e>1>s\tanh(s/2)$. Hence $-2$ is not a pole.

\subsection{Simplicity residues and the exact pole set}
At a genuine pole $\zeta$, the preceding argument gives
$Q(\zeta)\ne0$ and $\zeta\ne-2$. If its order is $m$, the most singular
orders in \eqref{eq:riccati} are $m+1$ from $P'$ and $2m$ from $P^2$;
the other terms are of lower order. Thus $m+1=2m$ and $m=1$.
Writing $P=a_{-1}/(z-\zeta)+O(1)$ and comparing coefficients gives
\eqref{eq:residue}.

Away from $Q=0$, a zero of $D$ cannot cancel against $N$. Indeed the two
linear combinations in \eqref{eq:trig} have coefficient determinant
\[
 (2E+1)^2+r^2=4E(2E+1)\ne0.
\]
The last inequality follows because $E\ne0$ and $2E+1=0$ would imply
$Q=0$. Simultaneous vanishing would force $\sin(zr/4)=\cos(zr/4)=0$,
which is impossible. This proves exactly \eqref{eq:poleset}.

For a genuine pole define
\begin{equation}\label{eq:Czeta}
 C_\zeta=-\frac{\Res(P,\zeta)}\zeta
 =\frac{2e^{-2\zeta}(4e^\zeta-1)}{\zeta(\zeta+2)}.
\end{equation}
There is no pole at zero. For every fixed $R>0$ with no pole on
$|z|=R$, there are finitely many poles inside, and subtraction of their
principal parts leaves a function analytic on a neighborhood of the
closed disk. Cauchy's bound therefore proves
\begin{equation}\label{eq:finitepole}
 \frac{p_n}{n!}=\sum_{|\zeta|<R}C_\zeta\zeta^{-n}+O_R(R^{-n}).
\end{equation}
This is a finite-radius expansion. It does not assert convergence of an
infinite sum over all poles, a numerical ordering of nonreal poles, or an
effective bound for the constant implicit in the remainder.

\section{Differential algebraicity}\label{sec:algebraic}
Use independent jet variables $p,u,v$ and an auxiliary variable $X$.
Define the following polynomials in $\mathbb Q[z,p,u,v]$:
\begin{align*}
 A&=(z+2)p^2,& B_1&=-8u-4p+4(z+2),& C_1&=2u+2p-(z+2),\\
 A_2&=p^2+2(z+2)pu+2A,&
 B_2&=-8v-4u+4+B_1,& C_2&=2v+2u-1.
\end{align*}
The Riccati equation states
$AX^2+B_1X+C_1=0$ at $X=e^z$, $p=P$, $u=P'$. Differentiating with
$X'=X$ gives $A_2X^2+B_2X+C_2=0$ at $v=P''$. Therefore
\begin{equation}\label{eq:resultant}
 \mathcal F(z,p,u,v)=\det
 \begin{pmatrix}
 A&B_1&C_1&0\\
 0&A&B_1&C_1\\
 A_2&B_2&C_2&0\\
 0&A_2&B_2&C_2
 \end{pmatrix}
\end{equation}
satisfies $\mathcal F(z,P,P',P'')=0$.
It is not the zero polynomial: at $(z,p,u,v)=(0,1,0,0)$, the two
quadratics are $2X(X+2)$ and $5X^2+8X-1$. Their resultant is
$2^2(-1)(3)=-12$. Thus $P$ is differentially algebraic of order at most
two over $\mathbb Q(z)$, despite being non-D-finite. The polynomial
condition can have extraneous solution germs; neither minimal order nor
uniqueness from this resultant is asserted. Earlier unsuccessful bounded
D-algebraic fitting experiments do not establish differential
transcendence and do not contradict this annihilator.

\section{Exact quadratic arithmetic enumeration}\label{sec:recurrence}
Let $q_n=4-\delta_{n,0}$ for $n\geq0$, with $q_{-1}=0$. For $n\geq0$ put
\begin{align*}
 E_n&=\sum_{k=0}^n\binom nkp_k,&
 S_n&=\sum_{k=0}^n\binom nkp_kp_{n-k},\\
 T_n&=\sum_{k=0}^n\binom nk2^{n-k}S_k,&
 G_n&=\sum_{k=0}^{n-1}\binom nkp_{k+1},
\end{align*}
where $T_{-1}=G_0=0$. Extracting exponential coefficients in
\eqref{eq:riccati} gives
\begin{equation}\label{eq:recurrence}
 \boxed{6p_{n+1}=2p_n-4E_n+2q_n+nq_{n-1}+2T_n+nT_{n-1}-8G_n},
 \qquad p_0=1.
\end{equation}
For example, the left side of \eqref{eq:riccati} contributes
$6p_{n+1}+8G_n$, while the three terms on the right contribute
$2p_n-4E_n$, $2q_n+nq_{n-1}$, and $2T_n+nT_{n-1}$.

All inputs at stage $n$ are already available from $p_0,\ldots,p_n$ and
previous $S_k,T_k$. Each finite sum has $O(n)$ terms. Computing through
$p_N$ thus takes $O(N^2)$ arithmetic operations. The binomial row can be
streamed from $\binom n0=1$ by
$\binom n{k+1}=\binom nk(n-k)/(k+1)$; powers of two can likewise be
updated exactly. Storing the three sequences uses $O(N)$ scalar entries.
Divisibility by six is a consequence of the EGF identity, and is also
checked explicitly by the companion at every stage.

This improves the total-count bounds of $O(N^4)$ arithmetic operations
and $O(N^3)$ scalar storage in~\cite[Section 2]{cgp}
\cite[Theorem 11]{abh}. The comparison is in the constant-cost arithmetic
model used there. Integers grow with $n$; no $O(N^2)$ bit-complexity,
wall-clock guarantee, optimality, or corresponding run-refined bound is
claimed.

The independent verification route expands \eqref{eq:H}--\eqref{eq:ND}
as rational ordinary power series, divides by $D$, and multiplies the
$n$th coefficient by $n!$. It does not use the Riccati recurrence.
Both routes match all 71 fixture counts $p_0,\ldots,p_{70}$ exactly.
The fixture records separate provenance: $p_1,\ldots,p_{25}$ are the
posted OEIS values checked in the earlier source work; the longer
fixture was computed independently from the published endpoint
recurrence. Agreement is finite evidence, not the proof of the formula.

\section{Unique dominant pole and leading asymptotic}\label{sec:dominance}
For completeness we supply the positive-operator isolation argument,
rather than infer a complex gap from the positive pole equation.
Let $\mathcal K$ denote the integral operator with kernel $K$ from
\eqref{eq:K}. For real $t>0$, multiplication by $\sqrt{w_t}$ is boundedly
invertible on $L^2(X,\mu)$. Thus $A_t$ is similar to the compact
self-adjoint, positivity-preserving operator
\[
 S_t=M_{\sqrt{w_t}}\mathcal K M_{\sqrt{w_t}}.
\]
It need not be positive semidefinite. For any states $s,s'$, every
interval $(a,b)$ with
\[
 0<a<\min(\lo(s),\lo(s')),
 \qquad \max(\hi(s),\hi(s'))<b<1
\]
overlaps both. This set has positive area. Consequently $S_t^2$ has a
strictly positive kernel almost everywhere.

\begin{lemma}\label{lem:perron}
For such an operator, $\lambda(t)=\spr(A_t)=\|S_t\|>0$ is a simple
positive eigenvalue with positive left and right eigenfunctions.
All other eigenvalues have smaller modulus.
\end{lemma}
\begin{proof}
The compact self-adjoint spectral theorem gives a top eigenvector of
$S_t^2$. Replacing it by its absolute value cannot decrease its Rayleigh
quotient, so the absolute value is also a top eigenvector. Strict
positivity of the squared kernel makes it positive almost everywhere.
A sign-changing top eigenvector would yield a strictly larger Rayleigh
quotient after taking its absolute value, a contradiction. Hence the top
eigenspace is one-dimensional: two independent real vectors would
produce a nonzero vector orthogonal to a positive vector, which must
change sign. The action of $S_t$ on the positive vector is positive,
so its eigenvalue is $+\|S_t\|$. Any other eigenvalue of the same modulus
would give another top vector of $S_t^2$. Compactness gives the strict
spectral gap. Similarity transfers simplicity and positive left and
right eigenvectors to $A_t$.
\end{proof}
Norm continuity gives continuity of $\lambda$. Increasing $t$ strictly
increases the kernel of $S_t$ wherever $K=1$; testing on its positive
Perron vector proves strict increase of $\lambda$. Also $\lambda(t)\to0$
as $t\downarrow0$. On
$F=\{(a,b):a<1/4,\ b>3/4\}$, all pairs overlap, $\mu(F)=1/16$, and
$w_t\geq t^2e^{t/2}$. Its indicator Rayleigh quotient gives
\[
 \lambda(t)\geq t^2e^{t/2}/16\longrightarrow\infty.
\]
Thus there is a unique $t_*>0$ with $\lambda(t_*)=1$.

The simple eigenvalue depends analytically on the parameter locally.
If $f_t>0$ is its real unit eigenvector for $S_t$, differentiation gives
\begin{equation}\label{eq:lambdaprime}
 \lambda'(t)=\lambda(t)\int_X\frac{w'_t}{w_t}f_t^2\dd\mu>0.
\end{equation}
One can obtain this directly from
$S_t'=\tfrac12(M_{w'_t/w_t}S_t+S_tM_{w'_t/w_t})$ and the normalized
simple-eigenvalue derivative formula. Local analyticity follows from a
Riesz projection around the isolated eigenvalue, choosing the positive
weight square root locally.

At $t_*$ put $h_*=M_{1/\sqrt{w_{t_*}}}f_{t_*}$ and
$\ell_*=M_{\sqrt{w_{t_*}}}f_{t_*}$. Their pairing is one. The local
resolvent singular part is the corresponding analytic rank-one
projection divided by $1-\lambda(z)$. Its scalar numerator in
\eqref{eq:resolvent} is
\[
 \left(\int_X\sqrt{w_{t_*}}f_{t_*}\dd\mu\right)^2>0.
\]
Equation \eqref{eq:lambdaprime} proves that this is a genuine simple
positive scalar pole. For $|z|<t_*$, the absolute kernel comparison
$|A_z^kf|\leq A_{|z|}^k|f|$ and the spectral-radius formula show
$\spr(A_z)\leq\lambda(|z|)<1$; hence there are no poles inside.
The scalar resolvent identity extends over this disk by its invertible
analytic resolvent. It follows that $t_*=\rho_0$ from
\eqref{eq:realpoles}. Write $\rho=\rho_0$.

For $|z|=\rho$, $z\ne\rho$, define
$B_z=\mathcal K M_{|w_z|}$. Its weights are positive and bounded away
from zero. The same squared-kernel argument gives a positive Perron
vector $h_z$ with eigenvalue $\beta=\spr(B_z)$. On the interval component,
\[
 |w_z(a,b)|=\rho^2e^{\operatorname{Re}(z)(b-a)}<w_\rho(a,b),
\]
whereas the singleton weights have equal absolute value. Using the
positive left eigenvector $\ell_*$ of $A_\rho$ yields
\begin{align*}
 (1-\beta)\int_X\ell_*h_z\dd\mu
 &=\int_\Delta(w_\rho-|w_z|)h_z(\mathcal K\ell_*)\dd\mu>0.
\end{align*}
Indeed $\mathcal K\ell_*=\ell_*/w_\rho>0$ almost everywhere.
Thus $\spr(A_z)\leq\beta<1$ at every other point of this circle.
Invertibility near such a point and continuation from the interior
identify the quotient with the analytic scalar resolvent there.

We already know from \eqref{eq:entire} that scalar poles form a discrete
set. Compactness of the circle and discreteness, including at $\rho$,
give some $R>\rho$ with no other scalar pole on or inside that disk.
Subtracting the $\rho$ pole and applying Cauchy's estimate proves
\begin{theorem}\label{thm:leading}
There exist $R>\rho$ and $0<\theta=\rho/R<1$ such that
\begin{equation}\label{eq:leading}
 p_n=C\,n!\rho^{-n}\bigl(1+O(\theta^n)\bigr),\qquad
 C=\frac{2e^{-2\rho}(4e^\rho-1)}{\rho(\rho+2)}>0.
\end{equation}
The positive pole $\rho$ is the unique dominant pole.
\end{theorem}
The residue formula gives $C$ directly; equivalently
$C=e^{-\rho}r(\rho)/(\rho\phi'(\rho))$.
This proof establishes existence of a complex pole-free gap, but supplies
no certified numerical value of $R$, error constant, or onset index.
The real interval certificate below does not fill that separate gap.

\section{Rational certificates for the constants}\label{sec:certificate}
Let the terminating decimals below denote exact rationals:
\begin{align*}
 L&=1.11343904173672704376166152691808324014139016583344946615,\\
 U&=1.11343904173672704376166152691808324014139016583344946616,\\
 C_L&=0.695688549070635767995703168724110156574198350721,\\
 C_U&=0.695688549070635767995703168724110156574198350722.
\end{align*}
The companion proves with integer and rational arithmetic that
\begin{equation}\label{eq:brackets}
 L<\rho<U,\qquad C_L<C<C_U.
\end{equation}
The following remainder bounds make the certificate reproducible without
assuming the correctness of numerical transcendental functions.

For $0\leq x<202$, let $T_{200}(x)=\sum_{j=0}^{200}x^j/j!$. Each ratio
between consecutive terms in the tail after $x^{201}/201!$ is at most
$x/202$. Hence
\begin{equation}\label{eq:expbound}
 T_{200}(x)\leq e^x\leq T_{200}(x)+
 \frac{x^{201}/201!}{1-x/202}.
\end{equation}
Monotonicity of $e^x$ extends these bounds to interval inputs.
For a nonnegative rational $x$, put $S=10^{120}$ and
$k=\lfloor\sqrt{\lfloor xS^2\rfloor}\rfloor$ using integer square root.
Checked integer squares prove
\[
 (k/S)^2\leq x<((k+1)/S)^2,
\]
which encloses $\sqrt x$. For $0\leq x\leq1$, the sine and cosine Taylor
terms decrease in magnitude. Truncating each alternating series at index
60 gives an upper bound, because the last included term is positive;
subtracting the next term's magnitude gives a lower bound. Sine is
increasing and cosine decreasing on this interval, so their endpoint
bounds propagate interval inputs rigorously.

For $x=L,U$ the program evaluates
\[
 (2e^{x/2}+1)\cos\!\left(\frac{x\sqrt{4e^x-1}}4\right)
 -\sqrt{4e^x-1}\sin\!\left(\frac{x\sqrt{4e^x-1}}4\right)
\]
by those rational interval operations. It verifies that both sine/cosine
arguments are in $(0,1)$ and that $D(L)>0>D(U)$. In this bracket
$0<zr/4<1$ and $0<\alpha<\pi/4$, so $\phi<1+\pi/4<\pi$.
Only the first positive tangent pole can be crossed. Continuity and
strict increase from \eqref{eq:phiprime} therefore isolate $\rho_0$.
Finally direct positive interval substitution into
\[
 C=\frac{2(4e^\rho-1)}{e^{2\rho}\rho(\rho+2)}
\]
proves the amplitude bracket. The machine-readable certificate records
outward decimal enclosures of both denominator signs and the propagated
amplitude interval. Decimal endpoints are suggested values only until
these rational inequalities have been verified; no binary floating
point or fitted coefficient is trusted by the certificate.

\section{Stirling expansions and safe inversion}\label{sec:inverse}
For sufficiently large integers, Theorem~\ref{thm:leading} gives
\begin{equation}\label{eq:logH}
 \log p_n=\mathcal H(n)+\epsilon_n,\quad
 \mathcal H(x)=\log\Gamma(x+1)-x\log\rho+\log C,\quad
 |\epsilon_n|\leq M\theta^n
\end{equation}
for some fixed $M>0$, $0<\theta<1$ and onset $n_0$.
For each fixed $m\geq1$, the ordinary Stirling remainder then gives
\begin{align}
 \log p_n={}&n\log n-(1+\log\rho)n+\tfrac12\log n
 +\log(C\sqrt{2\pi})\notag\\
 &+\sum_{j=1}^m\frac{B_{2j}}{2j(2j-1)n^{2j-1}}
 +O(n^{-2m-1}).\label{eq:stirling}
\end{align}
Exponentiating a sufficiently long finite truncation gives, for every
fixed $K\geq0$,
\[
 p_n=C\sqrt{2\pi n}\left(\frac n{e\rho}\right)^n
 \left(\sum_{k=0}^K\frac{s_k}{n^k}+O(n^{-K-1})\right),
\]
where $s_k$ are the formal coefficients of
\[
 \exp\left(\sum_{j\geq1}\frac{B_{2j}t^{2j-1}}{2j(2j-1)}\right)
 =1+\frac t{12}+\frac{t^2}{288}-\frac{139t^3}{51840}+\cdots.
\]
These are fixed-order Poincar\'e expansions; no convergence of the
infinite Bernoulli series is claimed. Also
$p_{n+1}/p_n=((n+1)/\rho)(1+O(\theta^n))$.

Define $N(y)=\min\{n\geq0:p_n\geq y\}$ for $y>0$. Appending the new
maximum extends the last ascending run and preserves the overlap
inequalities, giving an injection from size $n$ to size $n+1$.
Thus $(p_n)$ is nondecreasing; the ratio formula makes it eventually
strictly increasing and unbounded. For large $y$, let $x$ be the unique
large solution of $\mathcal H(x)=\log y$. This is the eventually
increasing branch because
$\mathcal H'(x)=\psi(x+1)-\log\rho\sim\log x$ and
$\mathcal H''(x)=\psi'(x+1)>0$.

\begin{proposition}[Rounding-safe threshold]
If \eqref{eq:logH} holds for all $n\geq n_0$ and
\[
 x-2\geq n_0,\qquad \mathcal H'(x-2)>0,\qquad
 \delta(x)=\frac{2M\theta^{x-2}}{\mathcal H'(x-2)}<1,
\]
then
\begin{equation}\label{eq:ceil}
 \ceil{x-\delta(x)}\leq N(y)\leq\ceil{x+\delta(x)}.
\end{equation}
\end{proposition}
\begin{proof}
Set $a=\ceil{x-\delta}-1$ and $b=\ceil{x+\delta}$. Both lie in
$[x-2,x+2]$. The increasing positive derivative of $\mathcal H$ gives
$\mathcal H(x)-\mathcal H(a)>2M\theta^{x-2}$ and
$\mathcal H(b)-\mathcal H(x)\geq2M\theta^{x-2}$.
The absolute errors at $a,b$ are at most $M\theta^{x-2}$.
Hence $p_a<y\leq p_b$, and monotonicity gives the bracket.
\end{proof}
If the two ceilings agree, they determine the exact threshold; otherwise
exact counting or stronger information is required. In particular one
cannot automatically round an asymptotic root, even if its error tends
to zero. Here $\delta(x)=O(\theta^x/\log x)$.
The certified $\rho,C$ intervals alone do not provide effective
$M,\theta,n_0$ for this implication.

For a convenient real starting estimate, write $L_y=\log y$ and
$u=W(L_y/(e\rho))$ on the positive real Lambert branch, so
$x_0=L_y/u$ satisfies $x_0\log(x_0/(e\rho))=L_y$.
One Newton correction in the Stirling expansion gives
\begin{equation}\label{eq:inverseinit}
 x=x_0-\frac{\tfrac12\log x_0+\log(C\sqrt{2\pi})}{u+1}
 +O\!\left(\frac1{x_0u}\right).
\end{equation}
Indeed the displayed correction is $O(1)$, the omitted residual is
$O(1/x_0)$, and the derivative is asymptotic to $u+1$.
This continuous inverse estimate is not an exact integer-ceiling rule.

\section{Source status and claim boundaries}\label{sec:sources}
The underlying overlap characterization is due to Asinowski, Banderier,
Billey, Hackl and Linusson~\cite{abbhl}, with the later development by
Asinowski, Banderier and Hackl~\cite{abh}. Claesson, Gu\dh mundsson and
Pantone~\cite{cgp} give polynomial-time counting and extensive numerical
experiments. Their author PDF, page 2, leaves a total-count generating
function or closed enumeration formula open. Section 3 distinguishes
unsuccessful bounded symbolic fits from experimental differential
approximants. In particular those fits do not prove non-D-finiteness or
differential transcendence. Their Section 3.3 predicts positive and
nonreal simple poles and the leading asymptotic numerically.

The authors' linked repository~\cite{repo} already contains estimates of
$\rho$, $C$, positive poles, and nonreal poles. For orientation, its first
positive estimates begin $1.113439\ldots$, $2.417184\ldots$,
$3.076673\ldots$; these decimal observations are prior art. Neither their
numerical discovery nor a claimed numerical complex-pole ordering is a
contribution of this report. The new statements here are the exact
formula, equations, residues, rational bounds, and proofs.

The 2021 paper~\cite{abh}, Section 2.4, records the exact asymptotic
problem and the infinitely-many-singularities/non-D-finiteness question
as unresolved. Its Theorem 10 is a multivariate ordinary functional
equation, distinct from the present one-variable explicit EGF.
Theorem 11 gives $O(N^4)$ arithmetic operations and $O(N^3)$ simultaneous
scalar allocations for total counts. The counting paper~\cite{cgp},
Section 2, gives the same operation and memory orders and explicitly
separates constant-cost arithmetic from its multiprecision implementation.
Neither comparison is a bit-complexity claim.

The final published 2023 version of Defant and Williams'
\emph{Semidistrim Lattices}~\cite{dw} was also checked, not only its
2021 preprint. Its Section 11.2 and Figure 12 connect the type-A image
count to A307030, while Remark 11.3 records later progress on neighboring
problems. The cited subsequent Choi--Sun work~\cite{cs} concerns other
Coxeter/Tamari/root-poset settings and does not provide the present
classical type-A total-count EGF. The current OEIS source entry~\cite{oeis}
was read as revision 63 dated 18 November 2021; its lack of a formula
field corroborates but cannot establish novelty.

A bounded source search on 3 October 2026 covered these primary papers,
the authors' live repository, the current OEIS entry, and exact or
near-exact searches combining pop-stacked permutations or A307030 with
generating function, closed form, exponential, D-finite, Riccati,
differential equation, and quadratic enumeration. No matching explicit
total-count EGF, proved non-D-finiteness, exact nonlinear ODE, or quadratic
arithmetic total-count recurrence was located. This is a bounded search
result, not proof of worldwide priority; scholarly author and citation
review remains appropriate before any submission. No publication,
author contact, or repository edit was undertaken.

\section{Companion and reproducible release}\label{sec:companion}
The source package contains a standard-library Python companion with
three bounded public modes: count through at most $N=400$, independently
verify quotient coefficients through at most $N=70$, and certify the
fixed real constants. It rejects invalid types, negative indices and
out-of-range requests before allocating a large computation. The public
computational API performs no writes; command-line JSON goes to standard
output, and errors go to standard error. Mathematical gates use explicit
exceptions and remain active under Python's optimization flag.

The quotient check is intentionally a separate, slower exact series
calculation. It is not covered by the recurrence's quadratic arithmetic
bound. The certificate uses exact fractions, Taylor bounds and integer
square roots as proved in Section~\ref{sec:certificate}. The data fixture
and output receipts distinguish source values from computed values.

PDF reproduction uses a fresh exclusive output directory, disabled TeX
shell escape, fixed metadata, a bounded stabilization loop, and settled
reference/layout checks. ZIP reproduction uses a fixed allowlist, SHA-256
verification, fixed member metadata, and exclusive no-clobber output.
The I/O helpers reject symbolic-link path components and pin opened
parent directories. This assumes trusted source files and trusted output
parent directories; it is not a sandbox for malicious TeX or a hostile
user who can replace files inside those directories.

Normal and optimized tests, two clean PDF builds, two ZIP creations,
and all-page rendered inspection are release checks. A deterministic
archive packages only the allowlisted report, source, fixtures,
certificates and reproduction scripts; temporary logs and page images
are excluded. Consult the package README for the exact commands and
scope of the final receipts.

\newpage
\begin{thebibliography}{9}
\bibitem{abbhl}
A. Asinowski, C. Banderier, S. Billey, B. Hackl and S. Linusson,
\emph{Pop-stack sorting and its image: Permutations with overlapping runs},
Acta Math. Univ. Comenianae 88(3) (2019), 395--402.
\href{https://www.iam.fmph.uniba.sk/amuc/ojs/index.php/amuc/article/download/1307/672/}{Publisher PDF}.
Theorem 1 and Sections 3--4.
\bibitem{cgp}
A. Claesson, B. \'A. Gu\dh mundsson and J. Pantone,
\emph{Counting pop-stacked permutations in polynomial time},
\href{https://arxiv.org/abs/1908.08910}{arXiv:1908.08910};
\href{https://doi.org/10.1080/10586458.2021.1926001}{journal DOI}.
\href{https://akc.is/papers/033-Counting-pop-stacked-permutations-in-polynomial-time.pdf}{Author PDF},
pp. 2--6 and Sections 3.1--3.3; page references use this 12-page PDF.
\bibitem{abh}
A. Asinowski, C. Banderier and B. Hackl,
\emph{Flip-sort and combinatorial aspects of pop-stack sorting},
Discrete Math. Theor. Comput. Sci. 22:2 (2021).
\href{https://doi.org/10.46298/dmtcs.6196}{DOI};
\href{https://dmtcs.episciences.org/7411/pdf}{publisher PDF}.
Theorems 3, 10 and 11, Section 2.4 and conclusion.
\bibitem{repo}
B. \'A. Gu\dh mundsson and collaborators,
\href{https://github.com/SuprDewd/pop-stacked-perms}{Pop-stacked permutations repository},
README and
\href{https://github.com/SuprDewd/pop-stacked-perms/blob/master/singularity_estimates.txt}{singularity estimates}.
Checked 3 October 2026; numerical estimates, not certified bounds.
\bibitem{dw}
C. Defant and N. Williams, \emph{Semidistrim Lattices},
Forum Math. Sigma 11 (2023), e50.
\href{https://doi.org/10.1017/fms.2023.46}{DOI}.
Published Section 11.2, pp. 29--30; Figure 12, p. 31;
Remark 11.3, p. 33. Preprint: \href{https://arxiv.org/abs/2111.08122}{arXiv:2111.08122}.
\bibitem{cs}
Y. Choi and N. Sun, \emph{The image of the pop operator on various lattices},
\href{https://arxiv.org/abs/2209.13695}{arXiv:2209.13695}.
Scope checked against the primary abstract and the published reference
in \cite{dw}.
\bibitem{oeis}
The OEIS Foundation, \href{https://oeis.org/A307030}{A307030},
\href{https://github.com/oeis/oeisdata/blob/main/seq/A307/A307030.seq}{official source entry}.
Revision 63, 18 November 2021; checked 3 October 2026.
\bibitem{ekp}
R. Ehrenborg, S. Kitaev and P. Perry,
\emph{A spectral approach to consecutive pattern-avoiding permutations},
J. Combin. 2(3) (2011), 305--353.
\href{https://arxiv.org/abs/1009.2119}{arXiv:1009.2119}.
Background on integral-operator enumeration; the specific mixed-run
construction and calculations are proved here.
\end{thebibliography}
\end{document}
